\documentclass{article}

% packages.tex
% All \usepackage lines live here.
% Loaded from main.tex with % packages.tex
% All \usepackage lines live here.
% Loaded from main.tex with % packages.tex
% All \usepackage lines live here.
% Loaded from main.tex with \input{packages}, before \input{macros}.
\usepackage[a4paper, total={7in, 10in}]{geometry} % Set page and margins

% --- Maths ---
\usepackage{amsmath}     % core maths environments (align, etc.)
\usepackage{amssymb}     % extra symbols (\mathbb, etc.)
\usepackage{bm}          % bold maths symbols (\bm)
\usepackage{bbm}         % \mathbbm{1} for the indicator function
\usepackage{mathrsfs}    % \mathscr script font

% --- Add other thesis packages below as you need them, for example: ---
\usepackage{graphicx}          % figures
\usepackage{booktabs}          % nicer tables
\usepackage{physics}
\usepackage[dvipsnames]{xcolor} % coloured text

% Load hyperref LAST (after almost everything else):
\usepackage{hyperref}
\usepackage[acronym]{glossaries-extra}
, before % macros.tex
% All custom commands live here.
% Loaded from main.tex with \input{macros}, after \input{packages}.
% (Requires the packages loaded in packages.tex: bbm, mathrsfs.)
 
% --- Operators and named quantities ---
\newcommand{\TN}{\mathsf{TN}}
\newcommand{\var}{\mathbb{V}\mathrm{ar}}   % variance
\newcommand{\cov}{\mathbb{C}\mathrm{ov}}   % covariance
\newcommand{\corr}{\mathbb{C}\mathrm{orr}} % correlation
\DeclareMathOperator*{\argmin}{argmin}
\DeclareMathOperator{\tr}{tr}
 
% --- Distribution / convergence relations ---
\newcommand{\simapprox}{\stackrel{\mathrm{approx}}{\sim}}
\newcommand{\simiid}{\stackrel{\mathrm{iid}}{\sim}}
\newcommand{\simiidt}{\sim_\mathrm{iid}}
\newcommand{\cas}{\stackrel{\mathrm{a.s.}}{\longrightarrow}} % converges almost surely
\newcommand{\cd}{\stackrel{\mathrm{d}}{\longrightarrow}}     % converges in distribution
\newcommand{\cp}{\stackrel{\mathbb{P}}{\longrightarrow}}     % converges in probability
 
% --- Font shorthands ---
\newcommand{\bb}[1]{\mathbb{#1}}       % blackboard bold
\newcommand{\m}[1]{\mathbf{#1}}        % upright bold (matrices)
\newcommand{\vv}[1]{\boldsymbol{#1}}   % bold italic (vectors)   [was \v: caron accent]
\newcommand{\mc}[1]{\mathcal{#1}}      % calligraphic            [was \c: cedilla accent]
\newcommand{\s}[1]{\mathscr{#1}}       % script
 
% --- Upright constants and differentials ---
\newcommand{\di}{\mathrm{d}}           % differential d
\newcommand{\e}{\mathrm{e}}            % Euler's number
\newcommand{\iu}{\mathrm{i}\,}         % imaginary unit          [was \i: dotless i]
\newcommand{\I}{\mathbbm{1}}           % indicator function
\newcommand{\gvn}{\,|\,}               % "given" bar for conditionals
 .
\usepackage[a4paper, total={7in, 10in}]{geometry} % Set page and margins

% --- Maths ---
\usepackage{amsmath}     % core maths environments (align, etc.)
\usepackage{amssymb}     % extra symbols (\mathbb, etc.)
\usepackage{bm}          % bold maths symbols (\bm)
\usepackage{bbm}         % \mathbbm{1} for the indicator function
\usepackage{mathrsfs}    % \mathscr script font

% --- Add other thesis packages below as you need them, for example: ---
\usepackage{graphicx}          % figures
\usepackage{booktabs}          % nicer tables
\usepackage{physics}
\usepackage[dvipsnames]{xcolor} % coloured text

% Load hyperref LAST (after almost everything else):
\usepackage{hyperref}
\usepackage[acronym]{glossaries-extra}
, before % macros.tex
% All custom commands live here.
% Loaded from main.tex with % macros.tex
% All custom commands live here.
% Loaded from main.tex with \input{macros}, after \input{packages}.
% (Requires the packages loaded in packages.tex: bbm, mathrsfs.)
 
% --- Operators and named quantities ---
\newcommand{\TN}{\mathsf{TN}}
\newcommand{\var}{\mathbb{V}\mathrm{ar}}   % variance
\newcommand{\cov}{\mathbb{C}\mathrm{ov}}   % covariance
\newcommand{\corr}{\mathbb{C}\mathrm{orr}} % correlation
\DeclareMathOperator*{\argmin}{argmin}
\DeclareMathOperator{\tr}{tr}
 
% --- Distribution / convergence relations ---
\newcommand{\simapprox}{\stackrel{\mathrm{approx}}{\sim}}
\newcommand{\simiid}{\stackrel{\mathrm{iid}}{\sim}}
\newcommand{\simiidt}{\sim_\mathrm{iid}}
\newcommand{\cas}{\stackrel{\mathrm{a.s.}}{\longrightarrow}} % converges almost surely
\newcommand{\cd}{\stackrel{\mathrm{d}}{\longrightarrow}}     % converges in distribution
\newcommand{\cp}{\stackrel{\mathbb{P}}{\longrightarrow}}     % converges in probability
 
% --- Font shorthands ---
\newcommand{\bb}[1]{\mathbb{#1}}       % blackboard bold
\newcommand{\m}[1]{\mathbf{#1}}        % upright bold (matrices)
\newcommand{\vv}[1]{\boldsymbol{#1}}   % bold italic (vectors)   [was \v: caron accent]
\newcommand{\mc}[1]{\mathcal{#1}}      % calligraphic            [was \c: cedilla accent]
\newcommand{\s}[1]{\mathscr{#1}}       % script
 
% --- Upright constants and differentials ---
\newcommand{\di}{\mathrm{d}}           % differential d
\newcommand{\e}{\mathrm{e}}            % Euler's number
\newcommand{\iu}{\mathrm{i}\,}         % imaginary unit          [was \i: dotless i]
\newcommand{\I}{\mathbbm{1}}           % indicator function
\newcommand{\gvn}{\,|\,}               % "given" bar for conditionals
 , after % packages.tex
% All \usepackage lines live here.
% Loaded from main.tex with \input{packages}, before \input{macros}.
\usepackage[a4paper, total={7in, 10in}]{geometry} % Set page and margins

% --- Maths ---
\usepackage{amsmath}     % core maths environments (align, etc.)
\usepackage{amssymb}     % extra symbols (\mathbb, etc.)
\usepackage{bm}          % bold maths symbols (\bm)
\usepackage{bbm}         % \mathbbm{1} for the indicator function
\usepackage{mathrsfs}    % \mathscr script font

% --- Add other thesis packages below as you need them, for example: ---
\usepackage{graphicx}          % figures
\usepackage{booktabs}          % nicer tables
\usepackage{physics}
\usepackage[dvipsnames]{xcolor} % coloured text

% Load hyperref LAST (after almost everything else):
\usepackage{hyperref}
\usepackage[acronym]{glossaries-extra}
.
% (Requires the packages loaded in packages.tex: bbm, mathrsfs.)
 
% --- Operators and named quantities ---
\newcommand{\TN}{\mathsf{TN}}
\newcommand{\var}{\mathbb{V}\mathrm{ar}}   % variance
\newcommand{\cov}{\mathbb{C}\mathrm{ov}}   % covariance
\newcommand{\corr}{\mathbb{C}\mathrm{orr}} % correlation
\DeclareMathOperator*{\argmin}{argmin}
\DeclareMathOperator{\tr}{tr}
 
% --- Distribution / convergence relations ---
\newcommand{\simapprox}{\stackrel{\mathrm{approx}}{\sim}}
\newcommand{\simiid}{\stackrel{\mathrm{iid}}{\sim}}
\newcommand{\simiidt}{\sim_\mathrm{iid}}
\newcommand{\cas}{\stackrel{\mathrm{a.s.}}{\longrightarrow}} % converges almost surely
\newcommand{\cd}{\stackrel{\mathrm{d}}{\longrightarrow}}     % converges in distribution
\newcommand{\cp}{\stackrel{\mathbb{P}}{\longrightarrow}}     % converges in probability
 
% --- Font shorthands ---
\newcommand{\bb}[1]{\mathbb{#1}}       % blackboard bold
\newcommand{\m}[1]{\mathbf{#1}}        % upright bold (matrices)
\newcommand{\vv}[1]{\boldsymbol{#1}}   % bold italic (vectors)   [was \v: caron accent]
\newcommand{\mc}[1]{\mathcal{#1}}      % calligraphic            [was \c: cedilla accent]
\newcommand{\s}[1]{\mathscr{#1}}       % script
 
% --- Upright constants and differentials ---
\newcommand{\di}{\mathrm{d}}           % differential d
\newcommand{\e}{\mathrm{e}}            % Euler's number
\newcommand{\iu}{\mathrm{i}\,}         % imaginary unit          [was \i: dotless i]
\newcommand{\I}{\mathbbm{1}}           % indicator function
\newcommand{\gvn}{\,|\,}               % "given" bar for conditionals
 .
\usepackage[a4paper, total={7in, 10in}]{geometry} % Set page and margins

% --- Maths ---
\usepackage{amsmath}     % core maths environments (align, etc.)
\usepackage{amssymb}     % extra symbols (\mathbb, etc.)
\usepackage{bm}          % bold maths symbols (\bm)
\usepackage{bbm}         % \mathbbm{1} for the indicator function
\usepackage{mathrsfs}    % \mathscr script font

% --- Add other thesis packages below as you need them, for example: ---
\usepackage{graphicx}          % figures
\usepackage{booktabs}          % nicer tables
\usepackage{physics}
\usepackage[dvipsnames]{xcolor} % coloured text

% Load hyperref LAST (after almost everything else):
\usepackage{hyperref}
\usepackage[acronym]{glossaries-extra}

% macros.tex
% All custom commands live here.
% Loaded from main.tex with % macros.tex
% All custom commands live here.
% Loaded from main.tex with % macros.tex
% All custom commands live here.
% Loaded from main.tex with \input{macros}, after \input{packages}.
% (Requires the packages loaded in packages.tex: bbm, mathrsfs.)
 
% --- Operators and named quantities ---
\newcommand{\TN}{\mathsf{TN}}
\newcommand{\var}{\mathbb{V}\mathrm{ar}}   % variance
\newcommand{\cov}{\mathbb{C}\mathrm{ov}}   % covariance
\newcommand{\corr}{\mathbb{C}\mathrm{orr}} % correlation
\DeclareMathOperator*{\argmin}{argmin}
\DeclareMathOperator{\tr}{tr}
 
% --- Distribution / convergence relations ---
\newcommand{\simapprox}{\stackrel{\mathrm{approx}}{\sim}}
\newcommand{\simiid}{\stackrel{\mathrm{iid}}{\sim}}
\newcommand{\simiidt}{\sim_\mathrm{iid}}
\newcommand{\cas}{\stackrel{\mathrm{a.s.}}{\longrightarrow}} % converges almost surely
\newcommand{\cd}{\stackrel{\mathrm{d}}{\longrightarrow}}     % converges in distribution
\newcommand{\cp}{\stackrel{\mathbb{P}}{\longrightarrow}}     % converges in probability
 
% --- Font shorthands ---
\newcommand{\bb}[1]{\mathbb{#1}}       % blackboard bold
\newcommand{\m}[1]{\mathbf{#1}}        % upright bold (matrices)
\newcommand{\vv}[1]{\boldsymbol{#1}}   % bold italic (vectors)   [was \v: caron accent]
\newcommand{\mc}[1]{\mathcal{#1}}      % calligraphic            [was \c: cedilla accent]
\newcommand{\s}[1]{\mathscr{#1}}       % script
 
% --- Upright constants and differentials ---
\newcommand{\di}{\mathrm{d}}           % differential d
\newcommand{\e}{\mathrm{e}}            % Euler's number
\newcommand{\iu}{\mathrm{i}\,}         % imaginary unit          [was \i: dotless i]
\newcommand{\I}{\mathbbm{1}}           % indicator function
\newcommand{\gvn}{\,|\,}               % "given" bar for conditionals
 , after % packages.tex
% All \usepackage lines live here.
% Loaded from main.tex with \input{packages}, before \input{macros}.
\usepackage[a4paper, total={7in, 10in}]{geometry} % Set page and margins

% --- Maths ---
\usepackage{amsmath}     % core maths environments (align, etc.)
\usepackage{amssymb}     % extra symbols (\mathbb, etc.)
\usepackage{bm}          % bold maths symbols (\bm)
\usepackage{bbm}         % \mathbbm{1} for the indicator function
\usepackage{mathrsfs}    % \mathscr script font

% --- Add other thesis packages below as you need them, for example: ---
\usepackage{graphicx}          % figures
\usepackage{booktabs}          % nicer tables
\usepackage{physics}
\usepackage[dvipsnames]{xcolor} % coloured text

% Load hyperref LAST (after almost everything else):
\usepackage{hyperref}
\usepackage[acronym]{glossaries-extra}
.
% (Requires the packages loaded in packages.tex: bbm, mathrsfs.)
 
% --- Operators and named quantities ---
\newcommand{\TN}{\mathsf{TN}}
\newcommand{\var}{\mathbb{V}\mathrm{ar}}   % variance
\newcommand{\cov}{\mathbb{C}\mathrm{ov}}   % covariance
\newcommand{\corr}{\mathbb{C}\mathrm{orr}} % correlation
\DeclareMathOperator*{\argmin}{argmin}
\DeclareMathOperator{\tr}{tr}
 
% --- Distribution / convergence relations ---
\newcommand{\simapprox}{\stackrel{\mathrm{approx}}{\sim}}
\newcommand{\simiid}{\stackrel{\mathrm{iid}}{\sim}}
\newcommand{\simiidt}{\sim_\mathrm{iid}}
\newcommand{\cas}{\stackrel{\mathrm{a.s.}}{\longrightarrow}} % converges almost surely
\newcommand{\cd}{\stackrel{\mathrm{d}}{\longrightarrow}}     % converges in distribution
\newcommand{\cp}{\stackrel{\mathbb{P}}{\longrightarrow}}     % converges in probability
 
% --- Font shorthands ---
\newcommand{\bb}[1]{\mathbb{#1}}       % blackboard bold
\newcommand{\m}[1]{\mathbf{#1}}        % upright bold (matrices)
\newcommand{\vv}[1]{\boldsymbol{#1}}   % bold italic (vectors)   [was \v: caron accent]
\newcommand{\mc}[1]{\mathcal{#1}}      % calligraphic            [was \c: cedilla accent]
\newcommand{\s}[1]{\mathscr{#1}}       % script
 
% --- Upright constants and differentials ---
\newcommand{\di}{\mathrm{d}}           % differential d
\newcommand{\e}{\mathrm{e}}            % Euler's number
\newcommand{\iu}{\mathrm{i}\,}         % imaginary unit          [was \i: dotless i]
\newcommand{\I}{\mathbbm{1}}           % indicator function
\newcommand{\gvn}{\,|\,}               % "given" bar for conditionals
 , after % packages.tex
% All \usepackage lines live here.
% Loaded from main.tex with % packages.tex
% All \usepackage lines live here.
% Loaded from main.tex with \input{packages}, before \input{macros}.
\usepackage[a4paper, total={7in, 10in}]{geometry} % Set page and margins

% --- Maths ---
\usepackage{amsmath}     % core maths environments (align, etc.)
\usepackage{amssymb}     % extra symbols (\mathbb, etc.)
\usepackage{bm}          % bold maths symbols (\bm)
\usepackage{bbm}         % \mathbbm{1} for the indicator function
\usepackage{mathrsfs}    % \mathscr script font

% --- Add other thesis packages below as you need them, for example: ---
\usepackage{graphicx}          % figures
\usepackage{booktabs}          % nicer tables
\usepackage{physics}
\usepackage[dvipsnames]{xcolor} % coloured text

% Load hyperref LAST (after almost everything else):
\usepackage{hyperref}
\usepackage[acronym]{glossaries-extra}
, before % macros.tex
% All custom commands live here.
% Loaded from main.tex with \input{macros}, after \input{packages}.
% (Requires the packages loaded in packages.tex: bbm, mathrsfs.)
 
% --- Operators and named quantities ---
\newcommand{\TN}{\mathsf{TN}}
\newcommand{\var}{\mathbb{V}\mathrm{ar}}   % variance
\newcommand{\cov}{\mathbb{C}\mathrm{ov}}   % covariance
\newcommand{\corr}{\mathbb{C}\mathrm{orr}} % correlation
\DeclareMathOperator*{\argmin}{argmin}
\DeclareMathOperator{\tr}{tr}
 
% --- Distribution / convergence relations ---
\newcommand{\simapprox}{\stackrel{\mathrm{approx}}{\sim}}
\newcommand{\simiid}{\stackrel{\mathrm{iid}}{\sim}}
\newcommand{\simiidt}{\sim_\mathrm{iid}}
\newcommand{\cas}{\stackrel{\mathrm{a.s.}}{\longrightarrow}} % converges almost surely
\newcommand{\cd}{\stackrel{\mathrm{d}}{\longrightarrow}}     % converges in distribution
\newcommand{\cp}{\stackrel{\mathbb{P}}{\longrightarrow}}     % converges in probability
 
% --- Font shorthands ---
\newcommand{\bb}[1]{\mathbb{#1}}       % blackboard bold
\newcommand{\m}[1]{\mathbf{#1}}        % upright bold (matrices)
\newcommand{\vv}[1]{\boldsymbol{#1}}   % bold italic (vectors)   [was \v: caron accent]
\newcommand{\mc}[1]{\mathcal{#1}}      % calligraphic            [was \c: cedilla accent]
\newcommand{\s}[1]{\mathscr{#1}}       % script
 
% --- Upright constants and differentials ---
\newcommand{\di}{\mathrm{d}}           % differential d
\newcommand{\e}{\mathrm{e}}            % Euler's number
\newcommand{\iu}{\mathrm{i}\,}         % imaginary unit          [was \i: dotless i]
\newcommand{\I}{\mathbbm{1}}           % indicator function
\newcommand{\gvn}{\,|\,}               % "given" bar for conditionals
 .
\usepackage[a4paper, total={7in, 10in}]{geometry} % Set page and margins

% --- Maths ---
\usepackage{amsmath}     % core maths environments (align, etc.)
\usepackage{amssymb}     % extra symbols (\mathbb, etc.)
\usepackage{bm}          % bold maths symbols (\bm)
\usepackage{bbm}         % \mathbbm{1} for the indicator function
\usepackage{mathrsfs}    % \mathscr script font

% --- Add other thesis packages below as you need them, for example: ---
\usepackage{graphicx}          % figures
\usepackage{booktabs}          % nicer tables
\usepackage{physics}
\usepackage[dvipsnames]{xcolor} % coloured text

% Load hyperref LAST (after almost everything else):
\usepackage{hyperref}
\usepackage[acronym]{glossaries-extra}
.
% (Requires the packages loaded in packages.tex: bbm, mathrsfs.)
 
% --- Operators and named quantities ---
\newcommand{\TN}{\mathsf{TN}}
\newcommand{\var}{\mathbb{V}\mathrm{ar}}   % variance
\newcommand{\cov}{\mathbb{C}\mathrm{ov}}   % covariance
\newcommand{\corr}{\mathbb{C}\mathrm{orr}} % correlation
\DeclareMathOperator*{\argmin}{argmin}
\DeclareMathOperator{\tr}{tr}
 
% --- Distribution / convergence relations ---
\newcommand{\simapprox}{\stackrel{\mathrm{approx}}{\sim}}
\newcommand{\simiid}{\stackrel{\mathrm{iid}}{\sim}}
\newcommand{\simiidt}{\sim_\mathrm{iid}}
\newcommand{\cas}{\stackrel{\mathrm{a.s.}}{\longrightarrow}} % converges almost surely
\newcommand{\cd}{\stackrel{\mathrm{d}}{\longrightarrow}}     % converges in distribution
\newcommand{\cp}{\stackrel{\mathbb{P}}{\longrightarrow}}     % converges in probability
 
% --- Font shorthands ---
\newcommand{\bb}[1]{\mathbb{#1}}       % blackboard bold
\newcommand{\m}[1]{\mathbf{#1}}        % upright bold (matrices)
\newcommand{\vv}[1]{\boldsymbol{#1}}   % bold italic (vectors)   [was \v: caron accent]
\newcommand{\mc}[1]{\mathcal{#1}}      % calligraphic            [was \c: cedilla accent]
\newcommand{\s}[1]{\mathscr{#1}}       % script
 
% --- Upright constants and differentials ---
\newcommand{\di}{\mathrm{d}}           % differential d
\newcommand{\e}{\mathrm{e}}            % Euler's number
\newcommand{\iu}{\mathrm{i}\,}         % imaginary unit          [was \i: dotless i]
\newcommand{\I}{\mathbbm{1}}           % indicator function
\newcommand{\gvn}{\,|\,}               % "given" bar for conditionals
 
% acronyms.tex
%   \newacronym{label}{short form}{long form}
%     label      = the key you type inside \gls{...} (never printed)
%     short form = the abbreviation shown on the page
%     long form  = the full expansion (capitalise it however you want it to read)

\setabbreviationstyle[acronym]{long-short}

\newacronym{cut}{CUT}{cell under test}
\newacronym{cpi}{CPI}{coherent pulse interval}
\newacronym{t/r}{T/R}{transmit/receive}
\newacronym{stap}{STAP}{Space-Time Adaptive Processing}
\newacronym{sinr}{SINR}{Signal-to-Interference-plus-Noise Ratio}
\newacronym{scr}{SCR}{Signal-to-Clutter Ratio}
\newacronym{wmf}{WMF}{Whitening Matched Filter}
\newacronym{amf}{AMF}{Adaptive Matched Filter}
\newacronym{namf}{NAMF}{Normalised Adaptive Matched Filter}
\newacronym{cnr}{CNR}{clutter-to-noise ratio}
\newacronym{okd}{OKD}{optimum K-distributed detector}
\newacronym{ten}{TEN}{texture estimation and normalisation}
\newacronym{rmb}{RMB}{Reed-Mallett-Brennan}
\newacronym{scm}{SCM}{sample covariance matrix}
\newacronym{nscm}{NSCM}{normalised sample covariance matrix}
\newacronym{ml}{ML}{maximum likelihood}
\newacronym{prf}{PRF}{pulse repetition frequency}
\newacronym{mle}{MLE}{maximum likelihood estimate}
\newacronym{lrt}{LRT}{likelihood ratio test}
\newacronym{glrt}{GLRT}{generalised likelihood ratio test}
\newacronym{np}{NP}{Neyman--Pearson}
\newacronym{roc}{ROC}{receiver operating characteristic}
\newacronym{cfar}{CFAR}{constant false-alarm rate}
\addtolength{\footnotesep}{4pt}
\lstdefinestyle{matlab}{language=Matlab, basicstyle=\ttfamily\footnotesize,
      numbers=left, numberstyle=\tiny, breaklines=true, frame=single,
      keywordstyle=\color{blue!70!black}, commentstyle=\color{green!40!black},
      stringstyle=\color{red!60!black}, showstringspaces=false}

\title{Thesis Workbook}
\author{z5360316}
\date{September 2026}

\begin{document}

\maketitle

\section{Radar and Signals Background}
\subsection{Introduction}
Effective target detection requires the suppression of clutter and interference. This is achieved by applying adaptive filters to the \gls{cut}.

\subsection{Multichannel Space-Time (STAP) Array Radar}
Multichannel space-time array radar refers to advanced radar systems that combine multiple spatial antenna channels with temporal pulse processing—widely known as \Gls{stap}—to simultaneously filter out environmental clutter, mitigate jammers, and detect moving targets.
\begin{itemize}
    \item In the usual (monostatic) radar, the same $N$ elements both transmit and receive, but not at the same time and not in the same way.
    \item Transmit and receive are time-multiplexed. In a pulsed radar, each element transmits the pulse, then a \gls{t/r} switch flips it over to receive mode to listen for the echo.
    \item On transmit, the elements usually act together as one aperture. They radiate coherently - the phases are set so the whole array's energy adds up in one direction, forming a single transmit beam pointed where you're looking. It is effectively using the $N$ elements as one big antenna to concentrate power.
    \item On receive, the elements are kept separate. Each element (or small subarray) has its own receiver chain and its own digitiser, producing $N$ independent channels. Those $N$ simultaneous outputs are the spatial samples - the $N$ numbers compared to measure angle and to do adaptive spatial / space-time filtering.
\end{itemize}

\subsection{Airborne Radar Signal Model}
Consider an airborne radar with an along-track side looking linear array comprising $N$ antennas, with inter-element spacing $d$, transmitting $M$ pulses in a \gls{cpi} at the pulse repetition frequency, $f_{PRF}$. Let the radar be travelling at velocity $v_p$.  

\begin{figure}[h]
    \centering
    \includegraphics[width=0.2\linewidth]{figs/airborne_radar_signal_model.png}
    \caption{Airborne Radar Signal Model}
    \label{fig:airborne_radar_signal_model}
\end{figure}

The echo of each pulse over \underline{fast-time} is sampled. A pulse travels out and back at the speed of light $c$. Thus, an echo arriving a time $t$ after transmission originates from a range $R=\frac{c t}{2}$, so each fast-time sample corresponds to a distinct range. These samples are the \underline{range gates} (or range cells), and the radar collects $N_r$ of them spanning the observed range window, as shown in Figure~\ref{fig:airborne_radar_signal_model}.

For each of the $N_r$ range gates we obtain a full $N\times M$ space-time grid - one per (antenna, pulse) pair. Stacking these across all range gates yields the $N\times M\times N_r$  \underline{data cube}. Detection is then carried out one range gate at a time: the gate currently being tested for a target is the \gls{cut}, and the detector slides this cell across all $N_r$ range gates.

\subsubsection{Dimension 1: Space ($N$ antennas) encodes angle}
A wave coming from angle $\theta$ hits each of the $N$ antennas one after another. Between each antenna in the linear array, there is an additional path length of $d\sin{\theta}$ between neighbouring elements. 

\begin{figure}[h]
    \centering
    \includegraphics[width=0.35\linewidth]{figs/path_diff.png}
    \caption{Path Difference Derivation}
    \label{fig:path_diff}
\end{figure}

For a narrow-band signal of wavelength $\lambda$, this path difference shows up as a phase shift of $2\pi (\frac{d}{\lambda})\sin{\theta}$ between adjacent antennae. This is since, for a wave periodic in space with period $\lambda$, one full period (one wavelength of travel) corresponds to $2\pi$ radians of phase.

We call this the \underline{normalised spatial frequency}: $$f_s = \frac{d}{\lambda}\sin{\theta.}$$

Hence, each step along the array adds a phase of $2\pi f_s$. This leads to the \underline{spatial steering vector}: $$\vb s_\text{sp} = \left[1\quad e^{j2\pi f_s}\quad\cdots\quad e^{j2\pi(N-1)f_s}\right]^T.$$

\subsubsection{Dimension 2: Slow-time ($M$ pulses) encodes Doppler}
Between pulses (spaced by the pulse repetition interval $T_r = \frac{1}{f_{PRF}}$), a
target moving with radial velocity $v_r$ changes its range by
$\Delta R = v_r T_r$. Since each pulse travels to the target and back,
the round-trip path changes by $2\Delta R$, which for a wavelength $\lambda$
appears as a phase shift from one pulse to the next of:
$$\Delta\phi = \frac{2\pi}{\lambda}\,(2\Delta R)
             = \frac{4\pi}{\lambda}v_r T_r = 2\pi f_d.$$
As in the spatial case, a constant per-pulse phase step produces a phase ramp
across the $M$ pulses; its rate is the \underline{normalised Doppler frequency}:
$$f_d = \frac{2 v_r T_r}{\lambda} = \frac{2 v_r}{\lambda f_{PRF}},$$
where the factor of two (absent in the spatial case) comes from the two-way
propagation. Each pulse therefore adds a phase of $2\pi f_d$, giving the
\underline{temporal steering vector}:
$$\vb s_\text{temp} = \left[1 \quad e^{j2\pi f_d} \quad \cdots \quad
   e^{j2\pi(M-1)f_d}\right]^T.$$

Note that only the radial movement of the target is considered, since any tangential motion does not contribute to a change in range and hence has no Doppler effect.\\

Together, the two steering vectors capture the target's phase progression along the two axes: $\vb s_\text{sp}$ across the $N$ antennas (set by the angle) and $\vb s_\text{temp}$ across the $M$ pulses (set by the Doppler).

\paragraph{Relating $M$ and $f_{PRF}$.}
The $M$ pulses are slow-time samples taken at rate $f_{PRF}$, spaced by the
pulse repetition interval $T_r = \frac{1}{f_{PRF}}$. The total \underline{dwell} (coherent
processing interval) is therefore:
$$T_{CPI} = M\,T_r = \frac{M}{f_{PRF}}.$$
These set two different quantities: $f_{PRF}$ is the slow-time sampling rate,
fixing the unambiguous Doppler span (Nyquist) and hence the maximum
unambiguous radial velocity $v_{\max} = \lambda f_{PRF}/4$; while $M$ fixes the
Doppler resolution $\Delta f_d = \frac{f_{PRF}}{M} = \frac{1}{T_{CPI}}$ (and the integration gain).

\subsubsection{Combining both dimensions: Space-time snapshot}
A target at a given (angle, Doppler) produces, at each (antenna, pulse) sample, a phase that is the sum of its spatial phase (from the angle) and its temporal phase (from the Doppler) with respect to an arbitrary reference origin. The two phases add and thus the complex factors multiply: the sample at antenna $n$, pulse $m$ carries $(\vb s_\text{sp})_n(\vb s_\text{temp})_m$.\\

Collecting over all $(n, m)$ gives the $N\times M$ \underline{spatio-temporal steering matrix} given by the outer-product: $$\vb s = \vb s_\text{sp} \vb s_\text{temp}^T.$$

\subsubsection{The Received Signal}
The received signal for a single range gate is an $N\times M$ slice of the data cube, $$\m X = \alpha\m S + \m C + \m N,$$ where $\m S=\vb s_\text{sp}\vb s_\text{temp}^T$ is the target spatio-temporal steering matrix, $\m C$ is the clutter return, $\m N$ is the thermal noise and $\alpha$ is the (complex) target amplitude.\\

Vectorising by stacking the columns gives the length $L=NM$ received vector $$\vb x = \operatorname{vec}(\m X) = \alpha\,\vb s + \vb c + \vb n,\qquad \vb s = \operatorname{vec}(\m S) = \vb s_\text{temp}\otimes\vb s_\text{sp},$$ where $\vb s$ is the space-time \underline{steering vector}, $\vb c = \operatorname{vec}(\m C)$ and $\vb n = \operatorname{vec}(\m N)$. Writing the clutter-plus-noise (interference) as $\vv{\varsigma}=\vb c+\vb n$, $$\vb x = \alpha\,\vb s + \vv{\varsigma}.$$

Vectorising lets us treat the joint space--time snapshot with standard vector
detection theory, where a single $L\times L$ covariance matrix captures all spatial,
temporal and space--time correlations at once.\\

The clutter vector $\vb c = \sqrt{\tau} \vb q$ where the \underline{speckle} $\vb q$ is complex Gaussian distributed with zero mean and covariance matrix $\m G$ and $\tau$ is a random variable called the \underline{texture}. Thus, given $\tau$, the received clutter is distributed as $\mathcal{CN}(\vb 0, \m R)$ where $\m R = \tau \m G$. For K-distributed clutter, the texture is drawn from the gamma distribution such that $\tau\sim\Gamma(a, b)$, where $a$ and $b$ are the shape and scale respectively.

\subsection{Clutter-Plus-Noise Covariance}
The second-order statistics of the interference are captured by the $L\times L$
covariance matrix \begin{align}
    \m R &= \bb E[(\vv\varsigma-\vv\mu)(\vv\varsigma-\vv\mu)^H] \nonumber\\
         &= \bb E[\vv\varsigma\vv\varsigma^H], \nonumber
\end{align} 
since the clutter-plus-noise is zero-mean ($\vv\mu=\vb 0$), as both the speckle $\vb q$ and thermal noise $\vb n$ are zero-mean.This covariance matrix holds every spatial, temporal and space-time correlation of the clutter-plus-noise, since $i$ and $j$ range over all $L$ samples. \\

Under a Gaussian interference model, the optimal detector \underline{whitens} the interference with $\vb R^{-1}$ and matches to the target, giving a filter $\vb w \propto \vb R^{-1}\vb s$ and the matched-filter statistic $|\vb s^H \vb R^{-1}\vb x|^2$. See Appendix~\ref{app:wmf}.\\

In practice $\vb R$ is unknown and must be estimated from data -- typically a sample covariance matrix formed from training range gates surrounding the \gls{cut}. The quality of this estimate governs detection performance, and estimating $\vb R$ accurately (especially in heterogeneous, non-Gaussian sea clutter) is the central problem addressed in this work.

\subsection{The Clutter Ridge and Brennan's Rule}
Because the platform moves at velocity $v_p$, a stationary clutter patch at angle $\theta$ has a non-zero relative radial velocity equal to the projection of $v_p$ onto the line of sight, $v_r(\theta)=v_p\sin\theta$. Its Doppler is therefore $$f_{d,\text{clutter}} = \frac{2 v_p \sin\theta}{\lambda}\,T_r \;\propto\; \sin\theta.$$ But the spatial frequency is also $f_s=(d/\lambda)\sin\theta\propto\sin\theta$, so clutter maps to a \underline{straight line} in the angle--Doppler
$(f_s,f_d)$ plane --- the \underline{clutter ridge} --- with slope
$$\beta = \frac{f_{d,\text{clutter}}}{f_s} = \frac{2 v_p T_r}{d}.$$
A moving target generally has an additional radial velocity of its own, so it
falls \emph{off} the ridge (Figure~\ref{fig:clutter_ridge}). This is precisely
why joint angle--Doppler (\Gls{stap}) processing is needed: clutter and a slow
target may share the same angle \emph{or} the same Doppler, but they separate in
the two-dimensional plane.

\begin{figure}[h]
    \centering
    \includegraphics[width=0.5\linewidth]{figs/clutter_ridge.png}
    \caption{The clutter ridge in the angle--Doppler plane. Platform motion gives clutter a Doppler $\propto\sin\theta$, the same dependence as the spatial frequency, so clutter lies along a line of slope $\beta=2v_pT_r/d$. A target at the same angle but a different radial velocity falls off the ridge, and is separable only by joint space--time processing.}
    \label{fig:clutter_ridge}
\end{figure}

Because the clutter is confined to this ridge, it occupies a low-dimensional
subspace of the $L=NM$ space--time space. Brennan's rule estimates the clutter
rank as
$$r \approx N + \beta(M-1) \ll NM,$$
which is the basis for reduced-rank clutter suppression.

\newpage
\section{Detection and Statistical Theory}
\subsection{Hypothesis Test $H_0/H_1$}
The target detection problem is generally formulated as a hypothesis test with the null hypothesis, $H_0$ being that the target is not present and the alternative hypothesis $H_1$ being that the target is present. The following test statistic $y$ is formed and compared to a suitable threshold, $\gamma$,

\newpage
\section{Traditional Detectors}

\newpage
\section{Problem Statement and Research Gap}

\subsection{Motivation}
Detecting slow or small targets in sea clutter is difficult because sea clutter is
\underline{non-Gaussian} (heavy-tailed and ``spiky'') and \underline{heterogeneous}
(its statistics vary over range and time). The classical adaptive detectors --- the
adaptive matched filter (AMF), the normalised AMF (NAMF) and Kelly's generalised
likelihood-ratio test (GLRT) --- are all derived under a \underline{Gaussian,
homogeneous} clutter assumption. In realistic maritime scenes this assumption is
violated, so these detectors are mismatched and suffer a significant loss in
detection performance. A detector that correctly accounts for the compound-Gaussian
nature of sea clutter would improve maritime detection --- but, as explained below,
the \emph{optimum} such detector is not directly realisable.

\subsection{Background: The Compound-Gaussian Model}
Sea clutter is well modelled as a \underline{compound-Gaussian} process: the clutter
vector factors as $\vb c = \sqrt{\tau}\,\vb q$, where the \underline{speckle}
$\vb q\sim\mathcal{CN}(\vb 0,\m G)$ carries the fast correlation structure (through the
covariance $\m G$) and the \underline{texture} $\tau$ is a slowly-varying random power.
Conditioned on $\tau$, the interference covariance is $\m R = \tau\m G$; for
K-distributed clutter $\tau$ is Gamma-distributed. Detection performance hinges on this
interference covariance, through the optimal filter $\vb w\propto\m R^{-1}\vb s$.

\subsection{Current Literature}
\begin{itemize}
  \item \textbf{The optimum detector is clairvoyant.} Gini et al.\ (1998) derived the
  optimum coherent detector for K-distributed clutter (the OKD) from the
  likelihood-ratio test, but \emph{assuming the clutter statistics --- both $\m G$ and
  $\tau$ --- are known}. The authors note this is an upper bound on achievable
  performance, since a practical system must estimate these parameters. It is therefore
  a benchmark, not a usable detector.
  \item \textbf{The texture can be estimated.} Gini and Greco (2002) retrieve the
  texture $\tau$ from measured sea clutter using a cyclostationary (CM-RELAX) approach,
  and Aboutanios and Rosenberg (2020) estimate the per-cell texture and normalise it
  --- the texture-estimation-and-normalisation (TEN) detector --- to form a practical,
  adaptive detector.
  \item \textbf{But the speckle covariance is assumed, not estimated.} The TEN approach
  handles the texture while \emph{assuming the speckle covariance $\m G$ is constant
  (shared and homogeneous) across range cells}. This is convenient but only
  approximate: it fails when the clutter \emph{structure} --- not merely its power ---
  varies across the scene.
\end{itemize}

\subsection{The Research Gap}
Taken together, no practical scheme estimates \emph{both} compound-Gaussian covariance
components and feeds them into the optimum detector:
\begin{itemize}
  \item the optimum detector requires both $\tau$ and $\m G$, assumed known;
  \item existing adaptive methods estimate the \emph{texture} but \emph{assume or
  approximate the speckle covariance} $\m G$.
\end{itemize}
The missing piece is a \underline{robust estimate of the speckle covariance $\m G$}
for K-distributed clutter --- ideally one that relaxes the shared-$\m G$ assumption
--- which, combined with a texture estimate, would turn Gini's clairvoyant optimum
detector into a practical, adaptive one.

\subsection{Aim and Contribution}
This work develops an estimator of the speckle covariance $\m G$ for K-distributed sea
clutter, combines it with texture estimation, and uses the result in the Gini optimum
(OKD) detector. This realises a practical, adaptive form of the optimum
compound-Gaussian detector, whose detection performance is then evaluated against the
Gaussian-assuming baselines (AMF/NAMF).

\subsection{Scope}
Two levels of ambition are possible for the speckle-covariance estimate:
\begin{enumerate}
  \item estimate a single shared $\m G$ accurately, retaining the homogeneity
  assumption; or
  \item relax the shared-$\m G$ assumption, allowing $\m G$ to vary across the scene
  (e.g.\ in piecewise-homogeneous regions).
\end{enumerate}
This choice sets the risk and novelty of the project and is to be fixed in consultation
with the supervisor.


\newpage
\section{Appendix: Supporting Concepts}
\subsection{Optimal Linear Filter: The \Gls{wmf}}
\label{app:wmf}
We seek the linear filter $\vb w$ that best detects a target $\alpha\vb s$ in interference $\vv{\varsigma}$ of covariance $\vb R = \bb E[\vv{\varsigma}\,\vv{\varsigma}^H]$, from the received vector $\vb x = \alpha\vb s + \vv{\varsigma}$.

\paragraph{Output \Gls{sinr}.}
Passing $\vb x$ through $\vb w$ gives
\[
  y = \vb w^H\vb x
    = \underbrace{\alpha\,\vb w^H\vb s}_{\text{signal}}
    + \underbrace{\vb w^H\vv{\varsigma}}_{\text{interference}}.
\]
The signal power is $|\alpha|^2|\vb w^H\vb s|^2$. The interference power is 
\begin{align}
    \bb E\!\left[|\vb w^H\vv{\varsigma}|^2\right] &= \bb E\!\left[(\vb w^H\vv\varsigma)\cdot(\vb w^H\vv\varsigma)^H\right] \nonumber\\
    &= \bb E\!\left[(\vb w^H\vv\varsigma)\cdot(\vv\varsigma^H\vb w)\right] \nonumber\\
    &= \bb E\!\left[\vb w^H(\vv\varsigma\,\vv\varsigma^H)\,\vb w\right] \nonumber\\
    &= \vb w^H \bb E\!\left[\vv\varsigma\,\vv\varsigma^H\right] \vb w, \nonumber
\end{align}
since $\vb w$ is deterministic, and so 
$$\bb E\!\left[|\vb w^H\vv{\varsigma}|^2\right] = \vb w^H\vb R\,\vb w.$$ \\

The output \gls{sinr} is therefore
\[
  \mathrm{SINR}(\vb w) = \frac{|\alpha|^2\,|\vb w^H\vb s|^2}{\vb w^H\vb R\,\vb w}.
\]

\paragraph{Maximisation.}
Dropping the constant $|\alpha|^2$, we maximise $|\vb w^H\vb s|^2/(\vb w^H\vb R\,\vb w)$.
Factor $\vb R = \vb R^{1/2}\vb R^{1/2}$ and set $\vb u = \vb R^{1/2}\vb w$ and
$\vb v = \vb R^{-1/2}\vb s$, so that $\vb w^H\vb s = \vb u^H\vb v$ and
$\vb w^H\vb R\,\vb w = \vb u^H\vb u$. They simplify as such since $\m R$ is a covariance matrix and hence it is Hermitian (self-adjoint) and positive-definite (the thermal noise covariance $\sigma^2\m I$ keeps all eigenvalues $> 0$ \footnote{Thermal noise is modelled as $\vb n\sim\mathcal{CN}(\vb 0, \sigma^2\m I)$ - zero-mean and covariance $\sigma^2\m I$. This is because thermal noise comes from the receiver electronics and it has two defining properties. First, it is white noise (uncorrelated across samples), and so each channel's noise (each antenna's receiver, each pulse) is independent of others. Hence, all the off-diagonal covariance entries are zero. Second, there is equal power in every channel. The receivers are identical, so each sample has the same noise variance $\sigma^2$. The result is $\sigma^2 \m I$ has eigenvalues $>0$.}), so it has a Hermitian square root $\m R^{1/2}$ and $\m R^{1/2}$ is invertible ($\m R^{-1/2}$ exists). \\

By the Cauchy--Schwarz inequality, which states $|\langle \vb x, \vb y\rangle| \leq \lVert \vb x\rVert\,\lVert \vb y\rVert$,
\[
  \frac{|\vb w^H\vb s|^2}{\vb w^H\vb R\,\vb w}
  = \frac{|\vb u^H\vb v|^2}{\vb u^H\vb u}
  \;\le\; \vb v^H\vb v = \vb s^H\vb R^{-1}\vb s,
\]
with equality iff $\vb u \propto \vb v$, i.e.\ $\vb R^{1/2}\vb w \propto \vb R^{-1/2}\vb s$.
Hence the optimal filter is
\[
  \boxed{\;\vb w \propto \vb R^{-1}\vb s\;}
  \qquad\text{achieving}\qquad
  \mathrm{SINR}_{\max} = \vb s^H\vb R^{-1}\vb s.
\]

\paragraph{Interpretation: whiten, then match.}
Writing $\vb w \propto \vb R^{-1/2}\big(\vb R^{-1/2}\vb s\big)$, the filter output is $\vb w^H\vb x = (\vb R^{-1/2}\vb s)^H(\vb R^{-1/2}\vb x)$: the factor $\vb R^{-1/2}$ \underline{whitens} the interference (transforming its covariance to the identity)\footnote{Using $\mathrm{Cov}(\m A\vv\varsigma)=\m A\,\mathrm{Cov}(\vv\varsigma)\,\m A^H$ with $\m A=\m R^{-1/2}$:
$\mathrm{Cov}(\m R^{-1/2}\vv\varsigma)=\m R^{-1/2}\m R\,\m R^{-1/2}
=\m R^{-1/2}(\m R^{1/2}\m R^{1/2})\m R^{-1/2}=\m I$, since $\m R^{-1/2}$ is Hermitian.}, after which a plain matched filter to the whitened target $\vb R^{-1/2}\vb s$ is optimal. Equivalently, in the eigenbasis of $\vb R$, the inverse $\vb R^{-1}$ strongly attenuates directions of large interference power (the clutter subspace) and passes directions of small, noise-level power --- it nulls the clutter before matching to the target.

\paragraph{Remarks.}
\begin{itemize}
  \item Only the \emph{direction} of $\vb w$ matters; any scalar cancels in the test
  statistic. The MVDR/Capon normalisation fixes it as
  $\vb w = \vb R^{-1}\vb s/(\vb s^H\vb R^{-1}\vb s)$ (unit gain toward $\vb s$,
  minimum output power).
  \item Under \emph{Gaussian} interference this whitening matched filter is also the
  likelihood-ratio (Neyman--Pearson) optimal detector. For non-Gaussian
  (compound-Gaussian) clutter it remains the optimal \emph{linear} filter given
  $\vb R$, but the globally optimal detector is the Gini-type test.
  \item In practice $\vb R$ is unknown and replaced by an estimate $\hat{\vb R}$,
  giving the adaptive matched filter $|\vb s^H\hat{\vb R}^{-1}\vb x|^2$.
\end{itemize}

\newpage
\section{Action}
\subsection{Questions}
\begin{itemize}
    \item Are we doing multichannel space-time framing (like Ingara, $L = MN$) or the single-channel temporal one (like IPIX)?
    \item In Aboutanios 2020, we assume the speckle covariance matrix $\m G$ is constant across cells? Do I take this assumption too and try to estimate this constant $\m G$ or do I try to find varying $\m G$, then still use the \gls{ten}. For example, if I segment the range of $\m G$'s into groups where they are approximately equal, to get the texture estimate, and then to the optimum detector. i.e. Shared-$\m G$ scope: do I (a) estimate a single shared $\m G$ accurately (keeping homogeneity), or (b) relax it to a spatially-varying / piecewise-homogeneous $\m G$?
    \item In Aboutanios 2020, it is K-distributed clutter only. So we would look at the \gls{okd} detector in the Gini 1998 paper? \textcolor{red}{Yes}
    \item Which speckle-covariance estimator(s) should I start from --- \gls{scm}, \gls{nscm}, fixed-point/approximate-\gls{ml} (Tyler, Conte--De Maio--Ricci), or a structured/regularised form?
    \item Data: simulated only, or do I get access to real sea-clutter data (Ingara / DST)? Any access constraints?
    \item Which baselines and metrics should I target (\gls{amf}, \gls{namf}, clairvoyant \gls{okd}; $P_d$ vs \gls{sinr}/\gls{scr} at a fixed $P_{fa}$)?
\end{itemize}

\subsection{To-Do}
\subsubsection{Confirm / Waiting:}
\begin{itemize}
    \item Receive Melvin's clutter-model paper from Elias (before coding the simulation).
\end{itemize}

\subsubsection{Reading:}
\begin{itemize}
    \item Work through Gini (1998), Gini \& Greco (2002), Aboutanios \& Rosenberg (2020) in depth; keep notes in this workbook.
    \item Read the speckle-covariance-estimation literature: fixed-point / approximate-\gls{ml} (Conte--De Maio--Ricci; Pascal et al.), Tyler's estimator, and regularised approaches (e.g.\ SCoRE).
    \item Read Melvin's \gls{stap} clutter-model paper once received.
\end{itemize}

\subsubsection{Code:}
\begin{itemize}
    \item Build a K-distributed clutter $+$ noise simulator producing the $N\times M\times N_r$ data cube (Melvin model), parameterised by \gls{cnr}, shape $a$, \gls{prf}, $v_p$, etc.
    \item Implement the Gaussian baselines: \gls{amf} and \gls{namf}.
    \item Implement the clairvoyant \gls{okd} (Gini) with known $\m G$ and $\tau$.
    \item Reproduce a known result (e.g.\ a $P_d$-vs-\gls{sinr} curve from Aboutanios 2020 or Gini 1998) as a correctness check.
    \item Implement candidate speckle-covariance estimators (\gls{scm}, \gls{nscm}, fixed-point) and test estimation accuracy on simulated clutter.
    \item Set up a Monte-Carlo evaluation harness ($P_{fa}$ calibration, $P_d$ curves).
\end{itemize}

\subsubsection{Writing:}
\begin{itemize}
    \item Keep the workbook updated (derivations, literature notes).
    \item Draft outline for the Thesis A preliminary report.
    \item \textcolor{red}{End of W5:} Submit short summary (50-100 words) about your thesis work, e.g. what is the problem? How are you going to solve it? What are the challenges? 
\end{itemize}

\paragraph{Short Summary.}
Traditional target detectors can perform poorly when clutter is non-Gaussian and heterogeneous, violating their underlying assumptions. The optimum detector for K-distributed clutter is known but assumes the clutter statistics (texture and speckle covariance) are separately available, which is not the case in practice. This thesis builds on existing texture-estimation methods to investigate whether a practical form of the optimum detector can be realised and to assess its performance against traditional detectors. The main challenge is that these statistics are not directly observable from data, so it must be established whether they can be estimated accurately enough for the detector to be practical.

\end{document}
